% ============================================================
% ЗАКЛЮЧЕНИЕ
% ============================================================
\conclusion

В данной работе исследована проблема холодного старта при развёртывании моделей машинного обучения в бессерверном инференсе на базе Kubernetes, а также разработана специализированная архитектура распределённого кэширования для её решения.

\predheading{6.1 Выполненные задачи} В ходе работы последовательно решены все поставленные задачи.

\textit{Задача 1: анализ причин холодного старта} Установлено, что рост размеров ML-моделей --- с сотен мегабайт у CNN-архитектур до сотен гигабайт у современных LLM --- превратил этап загрузки весов из второстепенной детали в доминирующую составляющую задержки холодного старта. Применение закона Литтла показало, что при типичных нагрузках задержка в 75~с при $\lambda = 5$~запр./с порождает очередь из 375 ожидающих запросов, что делает scale-to-zero практически неприменимым без специальной оптимизации.

\textit{Задача 2: исследование существующих решений} Проведённый сравнительный анализ Alluxio, JuiceFS, Dragonfly, CubeFS и Rook/Ceph показал, что ни одна из систем не удовлетворяет в полной мере специфике задачи: Alluxio избыточно сложен и требует коммерческой лицензии для enterprise-возможностей, JuiceFS требует объектного хранилища как обязательной зависимости, Dragonfly лишён POSIX-интерфейса. Это обосновало необходимость специализированной разработки.

\textit{Задача 3: промежуточное NFS-решение} Разработано промежуточное решение на базе NFS, устранившее проблему повторных загрузок из внешнего репозитория. Первый pod скачивает модель, остальные читают из общей директории. Разработаны алгоритмы координации с использованием \texttt{flock} и Kubernetes Lease. Однако NFS создал новый SPOF и bottleneck пропускной способности, что потребовало следующего этапа.

\textit{Задача 4: механизмы координации в Kubernetes} Исследованы файловые блокировки (\texttt{flock}), ресурс Lease из Coordination API и алгоритм Leader Election из библиотеки \texttt{client-go}. Установлено, что per-model Lease обеспечивает надёжную координацию первой загрузки модели без состояния гонки, однако не устраняет centralized-чтение через NFS.

\textit{Задача 5: архитектура распределённого chunk-кэша} Спроектирована трёхуровневая схема доступа: local cache~$\to$~peer cache~$\to$~remote fallback. Каждая worker-нода хранит чанки модели на локальном NVMe-диске. При запросе чанка агент проверяет локальное хранилище, затем обращается к Redis за списком пиров, затем загружает чанк из внешнего репозитория.

\textit{Задача 6: Redis как реестр метаданных} Обосновано применение Redis с TTL-based heartbeat в качестве централизованного реестра активных агентов и доступных чанков. Показано, что Redis обеспечивает $O(1)$-поиск пиров с задержкой менее 1~мс, что значительно быстрее gossip-конвергенции и надёжнее mDNS в overlay-сетях Kubernetes.

\textit{Задача 7: применение FUSE} Обоснован выбор FUSE для предоставления ML-runtime стандартного файлового интерфейса. Механизм VFS-перехватов позволяет реализовать lazy loading: runtime-контейнер начинает чтение файла модели немедленно, тогда как физическая загрузка чанков происходит по мере обращений. Время до первого байта составляет менее 1~с.

\textit{Задача 8: разделение UDS и TCP} Формализовано применение Unix Domain Socket для канала \texttt{storage-mounter}~$\to$~\texttt{storage-agent} (нулевые накладные расходы на сетевой стек, доступен только на одной ноде) и TCP для межнодового канала \texttt{storage-agent}~$\to$~\texttt{storage-agent} (совместимость с Kubernetes networking, поддержка mTLS).

\textit{Задача 9: реализация прототипа} Реализованы ключевые компоненты: \texttt{storage-mounter} с FUSE-файловой системой на базе \texttt{go-fuse}, локальным chunk-кэшем и интерфейсом к HuggingFace Hub; \texttt{storage-agent} с chunk-cache на файловой системе и fallback-загрузкой. Прототип верифицирует архитектурную концепцию и служит основой для производственной реализации.

\textit{Задача 10: методика оценки эффективности} Построена расчётная модель на базе AMAT, показывающая ускорение доступа к чанку в 30--50 раз при прогретом кластере. Проведён анализ пяти ключевых сценариев отказов с описанием механизмов graceful degradation. Выявлены ограничения и определены направления дальнейшего развития.

\predheading{6.2 Практическая ценность} Предложенная архитектура обеспечивает ряд практически значимых результатов. Во-первых, снижается внешний трафик к модельным хабам: при $N$~репликах вместо $N \times S_{\mathrm{model}}$ байт передаётся $S_{\mathrm{model}}$ --- однократная загрузка независимо от числа реплик. Для модели 15~ГБ и 10 реплик экономия составляет 135~ГБ внешнего трафика за одно burst-событие. Во-вторых, ускоряется повторный старт: попадание в локальный NVMe-кэш обеспечивает время загрузки модели $\approx$2~с против 75~с при загрузке из интернета. В-третьих, устраняется NFS как SPOF: распределённый кэш не имеет единого сервера, деградация затрагивает только ноды с отказами. В-четвёртых, сохраняется полная совместимость с любым ML-runtime без изменения его кода: FUSE предоставляет стандартный файловый интерфейс.

\predheading{6.3 Ограничения} Система наиболее эффективна для крупных моделей ($>$5~ГБ); для малых моделей накладные расходы на FUSE и Redis могут превышать выигрыш от кэширования. Первый запуск в незаполненном кластере требует полной загрузки из внешнего репозитория. Redis в базовой конфигурации остаётся точкой отказа для функции peer discovery (устраняется HA-конфигурацией). P2P-трафик между агентами требует инструментирования для корректной наблюдаемости.

\predheading{6.4 Направления дальнейшего развития} На основе анализа реализованного прототипа и выявленных ограничений определены следующие направления развития:

\begin{itemize}[leftmargin=*]
  \item Полная реализация Redis-реестра в коде \texttt{storage-agent}: регистрация чанков после загрузки, heartbeat-обновление TTL-записей, поиск пиров при cache miss.
  \item Реализация P2P-чтения между агентами по TCP с поддержкой параллельной загрузки чанков с нескольких пиров одновременно.
  \item Реализация UDS-транспорта для канала \texttt{storage-mounter}~$\to$~\texttt{storage-agent} с устранением накладных расходов сетевого стека.
  \item Добавление LRU/TTL-политики вытеснения чанков с учётом популярности моделей по закону Ципфа.
  \item Реализация prefetching: при обнаружении последовательного паттерна чтения агент заблаговременно загружает следующие чанки.
  \item Добавление контрольных сумм (SHA-256) для верификации целостности чанков при P2P-передаче.
  \item Поддержка mTLS между агентами для обеспечения аутентификации в production-окружении.
  \item Расширение источников данных: Ollama Hub, ModelScope, S3-совместимое хранилище.
  \item Проведение полноценного benchmark на production-кластере с реальными GPU-нагрузками и измерением $T_{\mathrm{model\_ready}}$ для актуальных LLM-моделей.
\end{itemize}

Полученные результаты подтверждают, что предложенный подход позволяет существенно сократить задержку холодного старта ML-инференса в Kubernetes, сохраняя преимущества serverless-масштабирования и обеспечивая отказоустойчивость без единой точки отказа.
